\section{Where to go next}
\label{sec:next}

\subsection{A heavier starter: DummyHippo}
When this tutorial feels too small, try \textbf{DummyHippo}: a deliberately
heavy multi-file project (nested chapters, TikZ/PGFPlots, denser math, and a
larger bibliography). Use it as a tougher starting point once you are
comfortable with the patterns above.

Download or clone the folder from:
\begin{center}
\url{https://github.com/martino-vic/undertwig/tree/main/DummyHippo}
\end{center}

To bring it into this workspace:
\begin{enumerate}
  \item Open that link and download the \texttt{DummyHippo} folder
        (clone the repository, or use GitHub's download / raw folder copy).
  \item In the app, click \textbf{Import new project}.
  \item Choose the local \texttt{DummyHippo} directory---the one that contains
        \texttt{main.tex} (not only an inner \texttt{chapters} folder).
  \item Compile with \textbf{Convert}. For citations, run
        \textbf{Update Bibliography}, then \textbf{Convert} again.
\end{enumerate}

\subsection{Bibliography wrap-up}
Classic references such as \cite{lamport94} and \cite{knuth84} are resolved from
\texttt{references.bib} via
\begin{verbatim}
\bibliographystyle{plain}
\bibliography{references}
\end{verbatim}
at the end of \texttt{main.tex}. Remember the usual loop: compile $\to$
bibliography tool $\to$ compile again.

\subsection{Thanks}
Thanks for reading through all of this---now open a blank section and start
writing your own material, or import DummyHippo if you want a tougher compile.
